\begin{lemma}[Elementary ordered union at a parallel extension]
Let $T$ be a positive-rung ETT satisfying maximum property, let its last
step be a PE, and write $T_n$, $v_n$, $\varphi_n$, and $D_n$ for the
terminal ladder tree, supporting vertex, generating coloring, and stored
colors.  Keep $G,k,e_0$ fixed but allow the certified input coloring to
vary.  Assume that for every positive-rung ETT $T'$ with arbitrary
certified input coloring, maximum property, and rung number $s<n$, the
following six assertions hold.
\begin{enumerate}
\item $V(T')$ is elementary under its generating coloring
$\varphi'_s$.
\item Under every proper $(T'_{s+1},D'_s,\varphi'_s)$-stable coloring
leaving exactly $e_0$ uncolored, $T'_{s+1}$ is interchangeable for any
two colors of which at least one is missing on it.
\item At every PE slot, stability descends to the prefix before the
support and to the ladder stage indexed by the support's first
appearance; stored pairs at distinct PE slots with distinct supports
and equal first-appearance indices are disjoint.
\item At a terminal PE, the stored two-color component through its
support meets the terminal ladder tree only at that support.
\item Every nonminimum tree endpoint of a defective boundary color
precedes each later PE support or SE extension vertex permitted by its
first-appearance index.
\item Every proper terminal stable coloring leaving exactly $e_0$
uncolored is a $\varphi'_s\bmod T'_s$ coloring.
\end{enumerate}

Let $\sigma_n$ be a proper $(T_n,D_n,\varphi_n)$-stable coloring leaving
exactly $e_0$ uncolored.  If $J_n$ is a closure of $T_n(v_n)$ under
$\sigma_n$ and $T_n\vee J_n$ is their ordered union, then
$V(T_n\vee J_n)$ is elementary under $\sigma_n$.
\end{lemma}
\begin{proof}
Put $n=r(T)$ and inspect the terminal PE constructor in
\noderef{TashkinovSeriesStep}.  It supplies the selected stable coloring,
the second stable coloring, the colors $\delta_n,\gamma_n$, the support
$v_n$, and the final two-color switch defining $\varphi_n$.

First truncate the certified series immediately before its terminal
step.  Restricting the series data and its certificates to the first
$n-1$ positive slots gives an ETT whose terminal tree is $T_n$ and whose
generating coloring is $\varphi_{n-1}$; the ordered-prefix and strictness
fields are inherited from \noderef{TashkinovSeries} and
\noderef{ExtendedTashkinovTree}.  Its maximum property is the
corresponding initial part of \noderef{MaximumProperty}.  If $n>1$, this
restriction has positive rung number and the lower-rung elementarity
hypothesis applies.  If $n=1$, it is the initial tuple and
\noderef{TechnicalTheoremBase} gives the same conclusion.  Hence in all
cases $V(T_n)$ is elementary under $\varphi_{n-1}$.

The two stable relations in the PE constructor preserve every
missing-color set on $T_n$ by \noderef{StableColoring}.  Before the final
switch, $V(T_n)$ is consequently elementary.  The PE component meets
$T_n$ only at $v_n$: this is current assertion (iv), obtained from
\noderef{ExtensionBase} using the displayed lower package.  The switch
therefore changes incident colors at no other tree vertex, and at
$v_n$ exchanges the missing color $\gamma_n$ with the present color
$\delta_n$.  Hence no missing color is duplicated on $T_n$ after the
switch.  Thus $V(T_n)$ is elementary under $\varphi_n$ and, by the
vertex clause of the assumed stable relation, under $\sigma_n$:
\[
             V(T_n)\text{ is elementary under }\sigma_n.       \tag{1}
\]

Let $j=m(v_n)$ be the first ladder index relevant to the support.
Current assertion (iii), namely \noderef{ExtensionRules} applied to the
same generalized lower package, descends the assumed terminal stability
to the prefix through $v_n$ and to $T_{j-1}$.  When $j>1$, the
restricted ETT ending at $T_{j-1}$ has positive rung number below $n$,
so the sixth lower-rung assertion says that $\sigma_n$ is a
$\varphi_{j-1}\bmod T_{j-1}$ coloring.  When $j=1$, the same conclusion
is obtained directly from the initial closure certificate: stability
copies its augmentations edge by edge, and the assumed closure $J_n$
supplies the new terminal closure.  This is the zero-rung construction
used in \noderef{TechnicalTheoremBase}.

The exact edge-control conclusion of
\noderef{HkuTreeColorsControlled} identifies the displayed colors on the
intervening ladder segment.  The base data \noderef{HierarchyBase} say
that $J_n$ is a closure of the same ordered prefix
$T_n(v_n)=T_j(v_n)$ under $\sigma_n$.  Thus the universal clause of
\noderef{ModColoring}, or the preceding zero-rung construction, gives an
actual corresponding ETT presentation with terminal tree $J_n$ above
$T_{j-1}$.  \noderef{MaximumPropertyTransport} gives that presentation
maximum property when its rung number is positive.  The first two
lower-rung assertions then apply; in the zero-rung case their two needed
conclusions instead come from \noderef{TechnicalTheoremBase}.  Together
with the closedness field of the \noderef{TashkinovClosure} certificate,
this gives
\[
 V(J_n)\text{ is elementary and closed under }\sigma_n,        \tag{2}
\]
and every pair of colors of which one is missing on $J_n$ is
$J_n$-interchangeable.

Suppose, for a contradiction, that the ordered union is not elementary.
By (1), (2), and the exact vertex-set equality in
\noderef{OrderedTreeUnion}, there are
\[
 u\in V(T_n)-V(J_n),\qquad v\in V(J_n)-V(T_n)
\]
and a color $\alpha$ missing at both $u$ and $v$.  The PE certificate
puts $\delta_n$ missing at $v_n\in V(T_n)\cap V(J_n)$; (2) therefore
gives $\alpha\ne\delta_n$.  The first conclusion of
\noderef{HkuParallelBoundaryStructure} says that the
$\gamma_n$-boundary is the singleton $\{f_n\}$.  Its conditional
conclusion, using the already established pre-step elementarity, would
make that boundary empty if $\gamma_n$ were missing on $T_n$; the PE
certificate has $\gamma_n\ne\delta_n$.  Hence $\gamma_n$ is not missing
on $T_n$, so $\alpha\ne\gamma_n$ as well.

We next produce the paper's vertex $w$.  The
\noderef{MaximumDefectiveChoice} certificate has at least two exact
$\delta_n$-boundary identities before the terminal switch.  Its second
stable relation preserves those identities and their colors.  The
switch changes the selected identity $f_n$ to $\gamma_n$ and fixes the
other $\delta_n$-boundary identities.  Choose one of the latter.
The unconditional delta-boundary conclusion of
\noderef{HkuParallelBoundaryStructure} gives an incident tree vertex
$w\prec_{T_n}v_n$.  Hence $w$ lies both in $T_n$ and in the supporting
prefix contained in $J_n$.

The numerical clause of the initial \noderef{KTriple} gives
$k\geq\Delta(G)+1$, while the degree of $w$ is at most $\Delta(G)$;
therefore some color $\beta$ is missing at $w$.  Since $J_n$ is
elementary, $\delta_n$ is missing at the distinct vertex $v_n$, and
$\gamma_n$ is missing nowhere on $T_n$, we have
$\beta\notin\{\delta_n,\gamma_n\}$.

Both $\alpha$ and $\beta$ are missing on $T_n$ and differ from
$\delta_n$.  The conditional conclusion of
\noderef{HkuParallelBoundaryStructure} therefore says that neither has
a boundary edge at $T_n$.  On $J_n$, closedness gives the same
boundary-free conclusion.  Interchangeability on $J_n$ implies that the
two nontrivial two-color components meeting $J_n$ at $v$ and $w$ are the
same.  They are nontrivial because elementarity makes $\beta$ present at
$v$ and $\alpha$ present at $w$.
\noderef{TwoColorComponentSpanningPath} identifies the common component
as a path; missingness and properness make $v,w$ its two degree-one ends.
The same argument using (1) gives a path with ends $u,w$ on $T_n$.

Connected components are disjoint by
\noderef{TwoColorComponentsDisjoint}, so the two paths sharing $w$ lie
in one component.  The distinct vertices $u,v$ both miss $\alpha$ and
$w$ misses $\beta$, so each vertex-node has at most one incident edge of
the other color by \noderef{MissingColors} and properness.  Thus all
three have degree at most one in the two-color incidence graph.  A connected
component containing three such distinct endpoints cannot be a path.
This contradiction proves that no common missing color at distinct
vertices of $T_n\vee J_n$ exists.  Therefore the union is elementary,
as required.
\end{proof}
